\section{Types of Varieties}

\subsection{Products of varieties}

The goal of this section is to construct the categorical product of varieties.
First we will construct the product of SWFs.

Recall that when constructing the colimits of SWFs we knew that the underlying topological space was the colimit of the spaces.
This is not true for limits; but we will see that the underlying set is the limit, but endowed with a different topology.

Let $(X,\O_X)$ and $(Y,\O_Y)$ be SWFs, then we construct a topology on $Z=X\times Y$ as follows.
For $U\subseteq X$ and $V\subseteq Y$ open, and $f\colon U\times V\to k$ of the form $f(x,y)=\sum_ig_i(x)h_i(y)$ for $g_i\in k[U]$
and $h_i\in k[V]$, define
$$ D_{U\times V}(f) = \set{(x,y)\in U\times V}[f(x,y)\neq0] . $$
This collection of $D_{U\times V}(f)$ forms a basis for our topology on $Z$.
Call such $f$ {\emphcolor basic regular} relative to $U,V$.

This allows us to define a structure sheaf on $Z$.
For $W\subseteq Z$ open, let $\O_Z(W)$ be the set of functions $\phi\colon W\to k$ such that there is a basic open cover
$W\subseteq\bigcup_iD_{U_i\times V_i}(f_i)$ and restricting $\phi$ to $W\cap D_{U_i\times V_i}(f_i)$ is of the form $\phi_i/f_i$
where $\phi_i$ is basic regular relative to $U,V$.

Notice that the projections $\pi_X\colon Z\to X$ and $\pi_Y\colon Z\to Y$ are morphisms of SWFs.
Indeed, $\pi_X^{-1}(D_U(f))=D_U(f)\times Y=D_{U\times Y}(f\cdot1)$ is open.
Since $D_U(f)$ forms a basis of open sets of $X$, we see that $\pi_X$ is continuous.
And for $f\colon U\to k$ regular, $\pi^*_X(f)(x,y)=f(x)$ is regular as well since $1\in k[Y]$.

Let us denote $(Z,\O_Z)$ by $(X,\O_X)\times(Y,\O_Y)$, or just $X\times Y$.

\blemm

    $(Z,\O_Z)$ is the categorical product of $(X,\O_X)$ and $(Y,\O_Y)$ in $\SWF_k$.

\elemm

\bproof

    For $\phi_X\colon T\to X$ and $\phi_Y\colon T\to Y$, in order for this to factor through $\phi\colon T\to Z$ we must have that
    $\phi=(\phi_X,\phi_Y)$ (since $Z$ is the product of sets).
    It remains to show that $\phi$ is a morphism.

    For every open $U\subseteq X,V\subseteq Y$, we have $\phi^{-1}(U\times V)=\phi_X^{-1}(U)\cap\phi^{-1}_Y(V)$ is open.
    And for every $f\colon U\times V\to k$ basic regular of the form $f(x,y)=\sum_ig_i(x)h_i(y)$, we have
    $$ \phi^*(f) = \sum_i\phi^*_X(g_i)\phi^*_Y(h_i) \colon \phi^{-1}(U\times V) \to k $$
    which is regular as $\phi_X,\phi_Y$ are morphisms and $k[\phi^{-1}(U\times V)]$ is an algebra.

    Thus we have that
    $$ \phi^{-1}(D_{U\times V}(f)) = D_{\phi^{-1}(U\times V)}(\phi^*(f)) \subseteq T $$
    is open, so $\phi$ is continuous.

    Now we need to show that $\phi^*$ induces a map of $k$-algebras.
    But this is immediate, as it maps basic regular maps to regular maps, and thus quotients of basic regular to regular maps.
    And the rest follows by the sheaf condition.
    \qed

\eproof

Since limits are functorial, if we have morphisms $f_i\colon X_i\to Y_i$ for $i=1,2$, the product $f_1\times f_2\colon X_1\times X_2\to Y_1\times Y_2$
is a morphism as well.

\blemm

    Let $(X,\O_X)$ and $(Y,\O_Y)$ be SWFs and $X'\subseteq X,Y'\subseteq Y$ be subsets.
    Then the following two SWFs on $X'\times Y'$ are the same
    $$ (X',\O_X|_{X'}) \times (Y',\O_Y|_{Y'}) = \bigl((X,\O_X)\times(Y,\O_Y)\bigr)|_{X'\times Y'} . $$

\elemm

\bproof

    Let $Z_1,Z_2$ denote the SWFs; we want to show $Z_1=Z_2$.
    We do this by showing that the identity of $X'\times Y'$ gives a morphism in both directions.

    In one direction, $\id\colon Z_1\to Z_2$ is a morphism since it is the restriction of the embedding $X'\times Y'\inj X\times Y$.
    This is just the product of the inclusions $(X',\O_X|_{X'})\to(X,\O_X)$ with the same one for $Y$, which is a morphism.

    Conversely we show by the universal property that $(X\times Y)|_{X'\times Y'}\to X$ obtained by projecting is a morphism.
    But this is just the restriction of the projection $X\times Y\to X$, which is a morphism.
    \qed

\eproof

\blemm

    For every two sets $X,Y$, there is a natural injective homomorphism
    $$ \mu\colon \fun(X,k) \otimes_k \fun(Y,k) \to \fun(X\times Y,k),\qquad \mu(f\otimes g)\colon(x,y)\mapsto f(x)g(y) . $$

\elemm

\bproof

    $\mu$ is clearly bilinear, so it is well-defined.
    Since $\fun(X,k),\fun(Y,k)$ are $k$-vector spaces, to show injectivity it is sufficient to show that $\mu$ maps a basis to a linearly
    independent set.
    Take a basis $\set{f_i}_{i\in I}$ of $\fun(X,k)$ and a basis $\set{g_j}_{j\in J}$ of $\fun(Y,k)$.
    Then $\set{f_i\otimes g_j}_{i\in I,j\in J}$ is a basis for their tensor product, and $\mu$ maps this to $f_ig_j$.
    Suppose we have a finite zeroing sum
    $$ \sum_{i,j}\alpha_{ij}f_ig_j = 0 . $$
    Then for every $y\in Y$ we have the zeroing sum
    $$ \sum_i\sum_j\alpha_{ij}g_j(y)f_i = 0 , $$
    which means that for all $i$, $\sum_j\alpha_{ij}g_j(y)=0$ by linear independence.
    Since this is true for all $y$ this means $\sum_j\alpha_{ij}g_j=0$, and thus $\alpha_{ij}=0$.
    \qed

\eproof

\bthrm

    If $X,Y$ are affine varieties, then $X\times Y$ is an affine variety.
    Furthermore, the natural map $k[X]\otimes_kk[Y]\to k[X\times Y]$, $f\otimes g\mapsto\pi_X^*(f)\cdot\pi_Y^*(g)$ is an isomorphism.

    If $X,Y$ are algebraic varieties, then so is $X\times Y$.

\ethrm

\bproof

    We first prove the first point, where $X,Y$ are affine.
    Since $k[X],k[Y]$ are fg $k$-algebras, so is $k[X]\otimes_kk[Y]$.
    Furthermore, $k[X]\subseteq\fun(X,k)$ and $k[Y]\subseteq\fun(Y,k)$, and as $k$ is a field all $k$-modules are free and thus flat so tensoring
    preserves subobjects, thus $k[X]\otimes_kk[Y]\subseteq\fun(X,k)\otimes_k\fun(Y,k)$.
    By the previous lemma this is a subalgebra of $\fun(X\times Y,k)$ which is reduced.
    Thus $k[X]\otimes_kk[Y]$ is a reduced fg $k$-algebra.

    This means there is an affine $k$-variety $Z$ such that $k[Z]=k[X]\otimes_kk[Y]$.
    Given any other SWF, we have natural bijections
    $$ \eqalign{
        \hom_\SWF(W,Z) &\cong \hom_\alg(k[Z],\O_W(W))\cr
            &\cong \hom_\alg(k[X]\otimes_kk[Y],\O_W(W))\cr
            &\cong \hom_\alg(k[X],\O_W(W)) \times \hom_\alg(k[Y],\O_W(W))\cr
            &\cong \hom_\SWF(W,X) \times \hom_\SWF(W,Y)
    } $$
    where the first is due to a result we showed in general: the global sections functor $\Gamma$ is a natural isomorphism for $Z$ affine.
    The third is because tensoring over $k$ is the coproduct of $k$-algebras.
    Thus $Z$ satisfies the universal property of products, so $Z\cong X\times Y$ is affine.

    Alternatively, we could note that $\bA^n\times\bA^m\cong\bA^{n+m}$ via the identity.
    Then for $X\subseteq\bA^n,Y\subseteq\bA^m$, $X\times Y\subseteq\bA^{n+m}$ is closed as it is equal to $Z(\I(X)\cup\I(Y))$.
    Thus it is affine.

    Now in general, let $X=\bigcup_iX_i$ and $Y=\bigcup_jY_j$ be open affine coverings.
    Then $X\times Y=\bigcup_{i,j}X_i\times Y_j$ is an open affine cover by the first point, so $X\times Y$ is algebraic.
    \qed

\eproof

Note that since $A\otimes_kk[x_1,\dots,x_n]\cong A[x_1,\dots,x_n]$ naturally, we see that
$$ k[X\times\bA^n] \cong k[X]\otimes_kk[\bA^n] \cong k[X][x_1,\dots,x_n] $$
for $X$ affine.

\bcoro

    Let $X,Y$ be algebraic varieties, $X'\subseteq X,Y'\subseteq Y$ be open/closed subvarieties.
    Then $X'\times Y'\subseteq X\times Y$ is an open/closed subvariety.

\ecoro

\bproof

    $X'\times Y'$ is open/closed in $X\times Y$ and $X'\times Y'$ is a variety.
    Furthermore we showed that the induced SWF structure on $X'\times Y'$ from $X\times Y$ is the same as the product of induced structures
    from $X,Y$.
    Thus $X'\times Y'$ is a subvariety of $X\times Y$.
    \qed

\eproof

\bprop

    Let $X,Y$ be algebraic varieties, then
    \benum
        \item $\dim(X\times Y)=\dim X+\dim Y$,
        \item If $X$ and $Y$ are irrreducible, so is $X\times Y$.
    \eenum

\eprop

\bproof

    First we show (2): if $X,Y$ are irreducible varieties, then so is $X\times Y$.
    Suppose $X\times Y=Z_1\cup Z_2$ for $Z_i$ closed, then define
    $$ X_i = \set{x\in X}[x\times Y\subseteq Z_i] = \bigcap_{y\in Y}\set{x\in X}[(x,y)\in Z_i]. $$
    Notice that $X=X_1\cup X_2$, and $X_i$ are closed as each $\set{x\in X}[(x,y)\in Z_i]$ is closed as the preimage of $Z_i$
    under $x\mapsto(x,y)$.
    Thus $X=X_i$ for some $i$, thus $X\times Y\subseteq Z_i$ and we have the desired result.

    Note that the map $x\mapsto(x,y)$ is continuous as its components are the identity and the constant map, both of which are continuous.
    So it is even a morphism of SWFs.

    Now we claim that if $X,Y$ are affine then $\dim(X\times Y)=\dim X+\dim Y$.
    Let $\dim X=n,\dim Y=m$, then there are finite surjective maps $X\to\bA^n,Y\to\bA^m$.
    The product map $X\times Y\to\bA^{n+m}$ is then finite surjective, and thus $\dim(X\times Y)=\dim(\bA^{n+m})=n+m$ as desired.

    Then in general for $X=\bigcup_iX_i,Y=\bigcup_jY_j$ open affine covers, $X\times Y=\bigcup_{i,j}X_i\times Y_j$.
    Then we have
    $$ \dim(X\times Y) = \max_{i,j}\dim(X_i\times Y_j) = \max_{i,j}(\dim X_i+\dim Y_j) = \max_i\dim X_i + \max_j\dim Y_j = \dim X + \dim Y $$
    as desired.
    \qed

\eproof

\bthrm

    The product $\bP^n\times\bP^m$ is projective.

\ethrm

\bproof

    Define
    $$ S\colon\bP^n\times\bP^m\to\bP^{(n+1)(m+1)-1},\qquad (\ell_1,\ell_2)\mapsto\ell_1\otimes\ell_2 . $$
    Explicitly, this maps lines $\ell_1\subseteq k^{n+1}$ and $\ell_2\subseteq k^{m+1}$ to $\ell_1\otimes_k\ell_2\subseteq k^{(n+1)(m+1)}$
    which we then project down to $\bP^{(n+1)(m+1)-1}$.
    Alternatively $S$ maps $[x_0:\cdots:x_n]$ and $[y_0:\cdots:y_m]$ to
    $$ [z_{ij}]_{0\leq i\leq n,0\leq j\leq m},\qquad z_{ij} = x_i\cdot y_j . $$

    We claim that $S$ is a closed embedding, so $\im S$ is closed and $S$ induces an isomorphism $S\colon\bP^n\times\bP^m\xto\sim\im S$.
    Being a morphism, and even a closed embedding, is local on the target.

    So take the standard affine open covering $D(Z_{ij})\subseteq\bP^{(n+1)(m+1)-1}$, then
    $$ S^{-1}D(Z_{ij}) = \set{(X,Y)}[X_iY_j\neq0] = D(X_i) \times D(Y_j) \subseteq \bP^n \times \bP^m . $$
    So we need to show that $D(X_i)\times D(Y_j)\to D(Z_{ij})$ is a closed embedding.

    Both the source and target are affine varieties, so this is a closed embedding iff $S^*$ is a surjective morphism of coordinate rings.
    Identifying $D(X_i)\cong\bA^n$ by $[X_0:\cdots:X_n]\mapsto(X_0/X_i,\dots,X_n/X_i)$.
    Let $z_{rs}=Z_{rs}/Z_{ij}$, $x_r=X_r/X_i$, and $y_s=Y_s/y_j$, so $k[D_{ij}]\cong k[z_{rs}]_{r,s}$.
    Then
    $$ S^*\colon k[\bA^{(n+1)(m+1)-1}] \to k[\bA^n] \otimes_k k[\bA^m] \cong k[x_0,\dots,x_n,y_0,\dots,y_m] $$
    is given by mapping $z_{rs}=Z_{rs}/Z_{ij}$ to $X_rY_s/X_iY_j=x_ry_s$.
    In particular $S^*(z_{rj})=x_ry_j=x_r$, and similar for $y$, so that $S^*$ is indeed surjective.
    \qed

\eproof

\bdefn

    The closed embedding $S$ in the above proof is called the {\emphcolor Segre embedding}.

\edefn

\bcoro

    The product of projective varieties is projective, and the product of quasi-projective/quasi-affine varieties is quasi-projective/quasi-affine.

\ecoro

\bproof

    For $X\subseteq\bP^n,Y\subseteq\bP^m$ closed, then $X\times Y$ is a closed subvariety of $\bP^n\times\bP^m$, and thus projective.
    The rest follows.
    \qed

\eproof

\subsection{Applications to dimension theory}

\bdefn

    For a variety $X$, consider the morphism $\Delta_X\colon X\to X\times X$ given by $\Delta_X(x)=(x,x)$.
    This is a morphism since its components are the identity; it is called the {\emphcolor diagonal morphism}.

\edefn

Note that the image $\Delta_X(X)=\set{(x,x)}[x\in X]$ is isomorphic to $X$ via $\Delta_X$.
Indeed, its inverse is the composition
$$ \Delta_X(X) \inj X\times X \xtos{\proj_1} X . $$
Thus for every $Y_1,Y_2\subseteq X$, the diagonal morphism forms an isomorphism
$$ Y_1\cap Y_2 \xtos\sim \Delta_X(Y_1\cap Y_2) = \Delta_X(X) \cap (Y_1\times Y_2) . $$
The left-hand side has structure induced from $X$, while the right-hand side's structure is induced from $X\times X$.

\bexam

    Let $X=\bA^n$, then $X\times X=\bA^{2n}$ and $\Delta(\bA^n)\subseteq\bA^{2n}$ is given by $Z(x_1-y_1,\dots,x_n-y_n)$.

\eexam

\bdefn

    A subset $C\subseteq\bA^{n+1}$ is called a {\emphcolor cone} if it is $k^\times$-invariant and $0\in C$.

\edefn

If we let $\pi\colon\bA^{n+1}-0\to\bP^n$ be the projection, we see that $C$ is a cone iff $C=\pi^{-1}(\pi(C-0))\cup0$.
In particular the maps $C\mapsto[C]=\pi(C-0)$ and $X\mapsto C(X)=\pi^{-1}(X)\cup0$ define a bijection between cones $C\subseteq\bA^{n+1}$
and subsets of $\bP^n$.

\blemm

    \benum
        \item For every closed $X\subseteq\bP^n$, its cone $C(X)\subseteq\bA^{n+1}$ is closed.
        \item If $X$ is irreducible, then so is $C(X)$.
        \item $\dim C(X)=\dim X+1$.
    \eenum

\elemm

\bproof

    For $X\subseteq\bP^n$ closed, $X=Z_{\bP^n}(f_1,\dots,f_m)$ for homogeneous $f_i$.
    Then $C(X)=Z_{\bA^{n+1}}(f_1,\dots,f_m)$ is closed.

    For (2), if $X=Z_{\bP^n}(f_1,\dots,f_m)$ is irreducible, then $(f_1,\dots,f_m)$ is prime and thus $C(X)=Z_{\bA^{n+1}}(f_1,\dots,f_m)$
    is irreducible.

    If $X=\bigcup_iX_i$ is a union of closed irreducible subsets, then $C(X)=\bigcup_iC(X_i)$ is a union of closed irreducible subsets.
    Then $\dim C(X)=\max\dim C(X_i)$, so we can assume that $X$ is closed irreducible.
    Since $\pi^{-1}X=C(X)-0$ is open dense in $C(X)$, it is sufficient to show that $\dim\pi^{-1}X=\dim X+1$.

    Consider the standard open affine covering $\bP^n=\bigcup_iU_i$, then $X=\bigcup_iX\cap U_i$ and $\pi^{-1}X=\bigcup\pi^{-1}(X\cap U_i)$.
    So it is sufficient to show
    $$ \dim\pi^{-1}(X\cap U_i) = \dim(X\cap U_i) + 1 . $$
    Notice that $\pi^{-1}U_i=D(x_i)$, and the map
    $$ (\pi,\proj_i)\colon\pi^{-1}U_i \to U_i\times(\bA^1-0) $$
    Is an isomorphism with inverse
    $$ s\colon([x_0:\cdots:x_n],t) \mapsto \parens{\frac{x_0}{x_i}t,\dots,\frac{x_n}{x_i}t} . $$
    These are both clearly inverses, and morphisms.

    Thus
    $$ \pi^{-1}U_i \cong U_i\times(\bA^1-0) , $$
    and this induces an isomorphism
    $$ \pi^{-1}(U_i\cap X) \cong (U_i\cap X)\times(\bA^1-0) . $$
    This has dimension $\dim(U_i\cap X)+1$ as desired.
    \qed

\eproof

\bthrm

    \benum
        \item Let $X,Y\subseteq\bA^n$ be closed irreducible subvarieties.
        Then every irreducible component of $X\cap Y$ has dimension at least $\dim X+\dim Y-n$.
        In particular if $X\cap Y$ is nonempty, then $\dim(X\cap Y)\geq\dim X+\dim Y-n$.

        \item Let $X,Y\subseteq\bP^n$ be closed irreducible subvarieties such that $\dim X+\dim Y\geq n$.
        Then $X\cap Y$ is nonempty and every irreducible component of $X\cap Y$ is of dimension at least $\dim X+\dim Y-n$.
    \eenum

\ethrm

\bproof

    For (1), there is an isomorphism
    $$ X\cap Y \xtos\sim (X\times Y)\cap\Delta(\bA^n) \subseteq \bA^{2n} . $$
    Since $X$ and $Y$ are irreducible, so is $X\times Y$ and $\dim(X\times Y)=\dim X+\dim Y$.
    Furthermore,
    $$ (X\times Y)\cap\Delta(\bA^n) = Z_{X\times Y}(x_1-y_1,\dots,x_n-y_n) \subseteq X\times Y $$
    is the zeroset of $n$ regular functions and thus the dimension of its irreducible components are each at least
    $\dim(X\times Y)-n=\dim X+\dim Y-n$.

    For (2), we know that $C(X),C(Y)$ are closed irreducible subsets of $\bA^{n+1}$.
    Moreover, $0\in C(X)\cap C(Y)$ so they have nonempty intersection.
    Thus from (1)
    $$ \multline{
        \dim C(X)\cap C(Y) \geq \dim C(X) + \dim C(Y) - (n+1) = (\dim X + 1) + (\dim Y + 1) - (n+1) =\cr
            \dim X + \dim Y - n + 1 \geq 1 .
    } $$
    In particular $C(X)\cap C(Y)$ is not finite, so there is an $x\in C(X)\cap C(Y)-0$ so $\pi(x)\in X\cap Y$, meaning $X\cap Y$
    is nonempty.

    Let $Z\subseteq X\cap Y$ be an irreducible component, then choose $U_i$ such that $Z\cap U_i\neq\emptyset$.
    Then $X\cap U_i,Y\cap U_i$ are closed irreducible subsets of $U_i\cong\bA^n$ and $Z\cap U_i$ is an irreducible component
    of their intersection.
    The result then follows from (1),
    $$ \dim Z = \dim(Z\cap U_i) \geq \dim(X\cap U_i) + \dim(Y\cap U_i) - n = \dim X + \dim Y - n . \qed $$

\eproof

\bcoro

    Let $f_1,\dots,f_m$ be homogeneous non-scalar polynomials in $k[x_0,\dots,x_n]$, with $m\leq n$.
    Then the closed set $Z=Z_{\bP^n}(f_1,\dots,f_m)$ is nonempty, and the dimension of each irreducible component of $Z$ is at least $n-m$.

\ecoro

\bproof

    We induct on $m$, for $m=1$ this is clear.
    Let $Z_m=Z_{\bP^n}(f_1,\dots,f_m)$ then inductively we have
    $$ Z_{m+1} = Z_m\cap Z_{\bP^n}(f_{m+1}) $$
    and $\dim Z_m+\dim Z_{\bP^n}(f_{m+1})\geq n-m+(n-1)=n+(n-(m+1))\geq n$ since $m+1\leq n$.
    Thus $Z_{m+1}$ is nonempty by the above theorem.
    And an irreducible component of $Z_{m+1}$ has dimension at least $\dim Z_m+\dim Z_{\bP^n}(f_{m+1})-n\geq n-(m+1)$ as desired.
    \qed

\eproof

\blemm

    Let $\phi\colon X\to Y$ be a dominant morphism of affine irreducible varieties.
    Then there exists an open dense subset $U\subseteq Y$ such that $\phi|_{\phi^{-1}U}\colon\phi^{-1}U\to U$ decomposes as
    $$ \phi = \phi^{-1}U \xtos g U\times\bA^n \xtos{\proj_U} U $$
    where $g$ is a finite surjective morphism.

\elemm

\bproof

    Since $\phi$ is dominant, $\phi^*\colon k[Y]\to k[X]$ is injective.
    Since $k[X]$ is a fg $k$-algebra, it is a fg $k[Y]$-algebra.
    Both $X,Y$ are irreducible and thus $k[X],k[Y]$ are domains so we can consider their fraction fields.
    Let $\p=0\subseteq k[Y]$, then the localization $k[X]_\p$ is a fg $k[Y]_\p$-algebra.

    By Noether normalization, there are $y_1,\dots,y_n\in k[X]_\p$ which are algebraically independent over $k(Y)$ and
    $k[X]_\p$ is finite over $k(Y)[y_1,\dots,y_n]$.

    We claim there is an $f\in k[Y]$ such that $y_1,\dots,y_n\in k[X]_f$ and $k[X]_f$ is a finite $k[Y]_f[y_1,\dots,y_n]$-algebra.

    Since $y_1,\dots,y_n\in k[X]_\p$ there is an $f'\in k[Y]$ such that $y_1,\dots,y_n\in k[X]_{f'}$ (take the product of the denominators
    of the $y_i$s).
    Let $a_1,\dots,a_r$ generate $k[X]$ as a $k[Y]$-algebra.
    Then $a_1,\dots,a_r$ are integral over $k(Y)[y_1,\dots,y_n]$ as $k[X]$ is finite over it, and thus integral over $k[Y]_{f''}[y_1,\dots,y_n]$
    for some $f''\in k[Y]$ as we can take common denominators.
    Then $f=f'f''$ does the job.

    With this in mind, the embeddings
    $$ k[Y]_f \inj k[Y]_f[y_1,\dots,y_n] \inj k[X]_f $$
    define morphisms $D_X(f)\xtos gD_Y(f)\times\bA^n\xtos\proj D_Y(f)$ of affine varieties.
    This is because these are all affine and
    $$ k[D_Y(f)\times\bA^n]\cong k[D_Y(f)]\otimes_kk[\bA^n]\cong k[Y]_f\otimes_kk[\bA^n] \cong k[Y]_f[y_1,\dots,y_n] . $$
    Since $k[Y]_f[y_1,\dots,y_n]\inj k[X]_f$ is injective and finite, $g$ is surjective and finite.
    \qed

\eproof

\bthrm

    Let $\phi\colon X\to Y$ be a dominant morphism of irreducible varieties.
    Then $f(X)$ contains a dense open subset $U\subseteq Y$ such that for every $y\in U$, $\phi^{-1}y$ is nonempty and each
    irreducible component of $\phi^{-1}y\subseteq X$ is of dimension $\dim X-\dim Y$.

\ethrm

\bnote

    Previously we showed that for any morphism of irreducible varieties, the fibers have dimension $\geq\dim X-\dim Y$.
    This theorem strengthens this for dominant morphisms.

\enote

\bproof

    For $X$ and $Y$ affine, then by the previous lemma there is a dense open $U\subseteq Y$ such that $\phi$ factors as
    $$ \phi^{-1}U \xtos g U\times\bA^n \xtos\proj U $$
    where $g$ is finite surjective.
    Since $U\subseteq Y$ and $\phi^{-1}U\subseteq X$ are open dense,
    $$ \dim X = \dim\phi^{-1}U = \dim U\times\bA^n = \dim U + n = \dim Y + n . $$
    So $n=\dim X-\dim Y$.

    And for $y\in U$, $g\colon\phi^{-1}U\to U\times\bA^n$ induces a finite surjective morphism $g\colon\phi^{-1}y\to\bA^n$.
    Thus $\dim\phi^{-1}y=n$, so $\phi^{-1}y\neq\emptyset$ and each irreducible component has dimension $\leq n=\dim X-\dim Y$.
    We know in general the dimension is at least $\dim X-\dim Y$ so we have equality.

    In general, we can replace $Y$ with a dense open affine subset $U\subseteq Y$ and $X$ by $\phi^{-1}U$.
    So we can assume that $Y$ is affine.
    Then choose an affine covering $X=\bigcup_iX_i$, then $\phi|_{X_i}\colon X_i\to Y$ is a dominant morphism of affine varieties
    (since $X_i$ is dense, $\phi(X_i)$ is dense in $\phi(X)$ which is then dense in $Y$).
    Thus each $\phi(X_i)$ contains an open dense $U_i$ such that for $y\in U_i$, the irreducible components of
    $\phi|_{X_i}^{-1}y=\phi^{-1}y\cap X_i$ all have dimension $\dim X_i-\dim Y=\dim X-\dim Y$.

    Let $U=\bigcap_iU_i$, then for each $y\in U$, $\phi^{-1}y=\bigcup_i\phi^{-1}y\cap X_i$ is an open cover.
    Since the irreducible components of $\phi^{-1}y\cap X_i$ have dimension $\dim X-\dim Y$, the same is true for $\phi^{-1}y$.
    \qed

\eproof

